\section{Introdução}
\label{sec:introducao}

Trocar o fundo de um vídeo é comum em videochamadas e transmissões ao
vivo: esconde o ambiente e protege a privacidade de quem está na câmera.
Para isso, o programa precisa decidir, em cada quadro e em tempo real,
quais pixels são da pessoa e quais são do fundo, um problema de
segmentação.

O ClearCut \citep{clearcut} faz isso com a webcam e troca o fundo com as
teclas A e D, mas segmenta com o \textit{SelfieSegmentation} do
MediaPipe, uma rede neural, o que não é permitido neste trabalho.
Mantivemos a experiência do ClearCut (webcam, pasta de fundos, teclas
A/D) e trocamos a segmentação por técnicas clássicas, com o conceito
novo da \textbf{subtração de fundo com MOG2}
\citep{zivkovic2004,opencv_tutorial_bgsub}.

O ClearCut-MOG2 funciona com a webcam ou com vídeos. Ele aprende a cena
vazia e, quando a pessoa entra, gera a máscara dela, conta os objetos
encontrados e cola a pessoa sobre o fundo escolhido, mostrando lado a
lado o original, a máscara e o resultado. Roda na CPU de um computador
comum, sem GPU nem modelo treinado; em troca, exige câmera parada e cena
vazia no início. Para avaliá-lo, gravamos dois vídeos com iluminação
diferente e comparamos as máscaras do programa com máscaras pintadas à
mão pelo grupo.
